\chapter{Compactness}
\label{ch:ii06}

Compactness is the finiteness principle of analysis. It converts infinite covers into finite subcovers, arbitrary sequences into convergent subsequences, and local continuity into global boundedness and uniform control.

\section*{Learning goals}
After this chapter, the reader should be able to:
\begin{itemize}
\item prove compactness using finite-subcover, sequential, or Heine--Borel criteria in the setting where each applies;
\item extract convergent subsequences from bounded sequences in Euclidean space;
\item use compactness to prove attainment of maxima and minima of continuous functions;
\item show that continuous images of compact sets are compact and use this to obtain boundedness or closedness conclusions;
\item construct open covers or sequences that demonstrate failure of compactness;
\item distinguish compactness from closedness and boundedness outside finite-dimensional Euclidean settings;
\end{itemize}
\section*{Conceptual roadmap}
\[
\boxed{\text{definition}\;\longrightarrow\;\text{estimate}\;\longrightarrow\;\text{structural theorem}\;\longrightarrow\;\text{application}.}
\]

\section{Open-cover compactness}
Every open cover of a compact space has a finite subcover.

\section{Sequential compactness}
In metric spaces compactness is equivalent to every sequence having a convergent subsequence with limit in the space.


% BEGIN VOL02-EXPANSION II06-example-01
\begin{example}[Sequential compactness extracts useful structure]\label{ex:ii06-expand-01}
Consider any sequence \((x_n)\) in the closed unit disk of \(\mathbb R^2\).
The disk is bounded, so Bolzano--Weierstrass supplies a convergent subsequence
\(x_{n_k}\to x\) in \(\mathbb R^2\). Because the disk is closed, \(x\) remains
in the disk. This verifies sequential compactness.

Both hypotheses in Heine--Borel appear explicitly: boundedness extracts a
subsequence, while closedness keeps its limit inside the set.
\end{example}
% END VOL02-EXPANSION II06-example-01
\section{Heine--Borel}
In Euclidean space compact sets are exactly the closed and bounded sets.

\section{Bolzano--Weierstrass}
Every bounded sequence in Euclidean space has a convergent subsequence.

\section{Continuous images}
Continuous images of compact spaces are compact.

\section{Extreme values}
Continuous real-valued functions on compact sets attain maxima and minima.


% BEGIN VOL02-EXPANSION II06-example-02
\begin{example}[Extreme values turn a supremum into an attained value]\label{ex:ii06-expand-02}
Let \(f(x)=x(1-x)\) on \([0,1]\). Compactness and continuity guarantee that a
maximum exists before we compute it. Completing the square,
\[
f(x)=\frac14-\left(x-\frac12\right)^2\le\frac14.
\]
Equality occurs at \(x=1/2\), so the maximum is \(1/4\).

The theorem supplies attainment; the algebra identifies the attaining point.
\end{example}
% END VOL02-EXPANSION II06-example-02
\section{Uniform continuity}
Continuous maps on compact metric domains are uniformly continuous.

\section{Completeness and total boundedness}
A metric space is compact exactly when it is complete and totally bounded.


% BEGIN VOL02-EXPANSION II06-example-03
\begin{example}[Total boundedness is stronger than boundedness]\label{ex:ii06-expand-03}
In an infinite-dimensional normed space, the closed unit ball is bounded but
need not be totally bounded. If it contained an infinite sequence of unit
vectors separated by a fixed positive distance, then no finite family of
sufficiently small balls could cover it.

This is why compactness is characterized by completeness plus total
boundedness, not merely completeness plus ordinary boundedness.
\end{example}
% END VOL02-EXPANSION II06-example-03
\section{Core structural results}

\begin{theorem}[Continuous image theorem]\label{thm:ii06-01}
If $K$ is compact and $f$ is continuous, then $f(K)$ is compact.
\begin{proof}
Pull an open cover of the image back to $K$, extract a finite subcover, and push the corresponding finite family forward.
\end{proof}
\end{theorem}

\begin{theorem}[Extreme-value theorem]\label{thm:ii06-02}
A continuous real function on nonempty compact $K$ attains its maximum and minimum.
\begin{proof}
Its image is compact in $\mathbb R$, hence closed and bounded; its supremum and infimum belong to the image.
\end{proof}
\end{theorem}

\begin{theorem}[Heine--Cantor]\label{thm:ii06-03}
A continuous map on a compact metric domain is uniformly continuous.
\begin{proof}
Otherwise choose pairs arbitrarily close with images separated by a fixed epsilon; a convergent subsequence of one side forces both sides to the same limit, contradicting continuity.
\end{proof}
\end{theorem}

\section{Worked examples}

\begin{example}[Open interval is not compact]\label{ex:ii06-worked-01}
Use an open cover to show $(0,1)$ is not compact. The sets $(1/n,1)$ for $n\ge2$ cover $(0,1)$ but no finite subfamily covers points near zero.
\end{example}

\begin{example}[Sequential noncompactness]\label{ex:ii06-worked-02}
Show $(0,1]$ is not compact using a sequence. The sequence $1/n$ has no subsequence converging to a point of $(0,1]$.
\end{example}

\begin{example}[Compact sets are closed]\label{ex:ii06-worked-03}
Prove this in metric spaces. For $x\notin K$, the continuous distance to $x$ attains a positive minimum on compact $K$, leaving a ball around $x$ disjoint from $K$.
\end{example}

\begin{example}[Compact sets are bounded]\label{ex:ii06-worked-04}
Prove it directly from an open cover. The increasing balls $B(x_0,n)$ cover the space; finitely many reduce to one largest ball.
\end{example}

\section{Solved dossiers}
Each dossier is a substantial solved problem. Corpus mappings guide topic choice where explicit retained source blocks exist; otherwise the dossier is a fresh canonical completion.

\begin{problem}[Open interval is not compact]\label{prob:ii06-dossier-01}
Use an open cover to show $(0,1)$ is not compact.
\end{problem}
\begin{solution}
The sets $(1/n,1)$ for $n\ge2$ cover $(0,1)$ but no finite subfamily covers points near zero.
\end{solution}

\begin{problem}[Sequential noncompactness]\label{prob:ii06-dossier-02}
Show $(0,1]$ is not compact using a sequence.
\end{problem}
\begin{solution}
The sequence $1/n$ has no subsequence converging to a point of $(0,1]$.
\end{solution}

\begin{problem}[Compact sets are closed]\label{prob:ii06-dossier-03}
Prove this in metric spaces.
\end{problem}
\begin{solution}
For $x\notin K$, the continuous distance to $x$ attains a positive minimum on compact $K$, leaving a ball around $x$ disjoint from $K$.
\end{solution}

\begin{problem}[Compact sets are bounded]\label{prob:ii06-dossier-04}
Prove it directly from an open cover.
\end{problem}
\begin{solution}
The increasing balls $B(x_0,n)$ cover the space; finitely many reduce to one largest ball.
\end{solution}

\begin{problem}[Continuous real functions are bounded]\label{prob:ii06-dossier-05}
Prove boundedness on compact domains.
\end{problem}
\begin{solution}
The compact image in $\mathbb R$ is bounded.
\end{solution}

\begin{problem}[Positive distance between disjoint compact sets]\label{prob:ii06-dossier-06}
Prove disjoint compact $K,L$ have positive distance.
\end{problem}
\begin{solution}
The continuous function $d(x,y)$ on compact $K\times L$ attains a minimum; zero would force a common point.
\end{solution}

\begin{problem}[Finite intersection property]\label{prob:ii06-dossier-07}
Prove compact spaces have it for closed sets.
\end{problem}
\begin{solution}
If the total closed intersection were empty, their complements would give an open cover with a finite subcover, contradicting the finite intersection hypothesis.
\end{solution}

\begin{problem}[Nested compact sets]\label{prob:ii06-dossier-08}
Show decreasing nonempty compact sets have nonempty intersection.
\end{problem}
\begin{solution}
They satisfy the finite intersection property inside the first compact set.
\end{solution}

\begin{problem}[Heine--Cantor contradiction sequence]\label{prob:ii06-dossier-09}
Describe the standard sequence proof.
\end{problem}
\begin{solution}
Failure gives $d(x_n,y_n)<1/n$ but image distance at least $\varepsilon_0$; compactness and continuity contradict this.
\end{solution}

\begin{problem}[Bolzano--Weierstrass by nested boxes]\label{prob:ii06-dossier-10}
Sketch the construction.
\end{problem}
\begin{solution}
Repeatedly choose a smaller closed box containing infinitely many sequence terms and diameters tending to zero; choose one term from each stage.
\end{solution}

\begin{problem}[Closed subset of compact is compact]\label{prob:ii06-dossier-11}
Prove it by covers.
\end{problem}
\begin{solution}
Add the open complement of the closed subset to any cover, extract a finite cover of the ambient compact set, then discard the complement.
\end{solution}

\begin{problem}[Compactness synthesis]\label{prob:ii06-dossier-12}
Why is compactness a finite-extraction principle?
\end{problem}
\begin{solution}
Its cover definition extracts finitely many neighborhoods, and its sequential form extracts one convergent subsequence from infinitely many terms.
\end{solution}

\section{Exercises with complete solutions}
The shorter exercise layer is kept separate from the dossiers.

\begin{exercise}\label{exr:ii06-01}
Is $\{1/n:n\ge1\}\cup\{0\}$ compact?
\end{exercise}
\begin{hint}
Check closed and bounded. In \(\mathbb R^n\), check closedness and boundedness first; Heine--Borel then converts them to compactness.
\end{hint}
\begin{solution}
Yes.
\end{solution}

\begin{exercise}\label{exr:ii06-02}
Is $\mathbb Z$ compact in $\mathbb R$?
\end{exercise}
\begin{hint}
Check boundedness. For sequential compactness, begin with an arbitrary sequence and construct a convergent subsequence.
\end{hint}
\begin{solution}
No.
\end{solution}

\begin{exercise}\label{exr:ii06-03}
Show every finite metric space is compact.
\end{exercise}
\begin{hint}
Choose one cover member per point. To show a continuous function attains a maximum, apply compactness to its image or to a maximizing sequence.
\end{hint}
\begin{solution}
Finitely many choices suffice.
\end{solution}

\begin{exercise}\label{exr:ii06-04}
Does compactness imply completeness in metric spaces?
\end{exercise}
\begin{hint}
Use a convergent subsequence of a Cauchy sequence. A continuous image of a compact set is compact; use this before trying to prove boundedness directly.
\end{hint}
\begin{solution}
Yes.
\end{solution}

\begin{exercise}\label{exr:ii06-05}
Is $(0,1)$ compact?
\end{exercise}
\begin{hint}
Use Heine--Borel. To disprove compactness, give either an open cover with no finite subcover or a sequence with no convergent subsequence.
\end{hint}
\begin{solution}
No.
\end{solution}

\begin{exercise}\label{exr:ii06-06}
Is a continuous image of compact closed in a metric codomain?
\end{exercise}
\begin{hint}
Compact subsets of metric spaces are closed. Closed subsets of compact spaces are compact; isolate the ambient compact set if one is available.
\end{hint}
\begin{solution}
Yes.
\end{solution}

\begin{exercise}\label{exr:ii06-07}
Why is $[a,b]$ compact?
\end{exercise}
\begin{hint}
Use Heine--Borel. Compact subsets of metric spaces are closed; separate a point outside the set using a positive distance argument.
\end{hint}
\begin{solution}
It is closed and bounded.
\end{solution}

\begin{exercise}\label{exr:ii06-08}
Must every sequence in a compact set converge?
\end{exercise}
\begin{hint}
Only a subsequence is guaranteed. Do not use Heine--Borel outside finite-dimensional Euclidean space without checking that the theorem still applies.
\end{hint}
\begin{solution}
No; $(-1)^n$ in $\{-1,1\}$ is a counterexample.
\end{solution}


% BEGIN VOL02-EXPANSION II06-exercises-01
\section{Graded supplementary exercises}
These exercises extend the preserved foundational set and are graded by mathematical role.

\subsection*{Standard computations}

\begin{exercise}[Heine--Borel test]\label{exr:ii06-expand-01}
Is \([0,1]\times[-2,3]\subset\mathbb R^2\) compact?
\end{exercise}
\begin{hint}
Check closedness and boundedness.
\end{hint}
\begin{solution}
Yes, by Heine--Borel.
\end{solution}

\begin{exercise}[Noncompact interval]\label{exr:ii06-expand-02}
Is \((0,1]\) compact in \(\mathbb R\)?
\end{exercise}
\begin{hint}
Check closedness or use the sequence \(1/n\).
\end{hint}
\begin{solution}
No; \(1/n\) converges to \(0\notin(0,1]\).
\end{solution}

\begin{exercise}[Continuous image]\label{exr:ii06-expand-03}
What can be said about \(f(K)\) if \(K\) is compact and \(f\) continuous?
\end{exercise}
\begin{hint}
Apply the continuous-image theorem.
\end{hint}
\begin{solution}
\(f(K)\) is compact.
\end{solution}

\begin{exercise}[Extreme value]\label{exr:ii06-expand-04}
Find the maximum of \(x^2\) on \([-2,1]\).
\end{exercise}
\begin{hint}
Compactness guarantees attainment; compare endpoint values.
\end{hint}
\begin{solution}
The maximum is \(4\) at \(x=-2\).
\end{solution}

\begin{exercise}[Finite epsilon-net]\label{exr:ii06-expand-05}
Give a \(1/2\)-net for \([0,1]\).
\end{exercise}
\begin{hint}
Choose finitely many evenly spaced centers.
\end{hint}
\begin{solution}
For example \(\{0,1/2,1\}\) covers \([0,1]\) by radius-\(1/2\) balls.
\end{solution}

\subsection*{Proofs}

\begin{exercise}[Continuous image compact]\label{exr:ii06-expand-06}
Prove continuous images of compact spaces are compact.
\end{exercise}
\begin{hint}
Pull an open cover back to the domain.
\end{hint}
\begin{solution}
The pullback covers \(K\); choose a finite subcover and push the corresponding finite family forward.
\end{solution}

\begin{exercise}[Compact metric sets bounded]\label{exr:ii06-expand-07}
Prove a compact metric space is bounded.
\end{exercise}
\begin{hint}
Cover it by balls \(B(x_0,n)\), \(n\in\mathbb N\).
\end{hint}
\begin{solution}
A finite subcover is contained in the largest selected ball.
\end{solution}

\begin{exercise}[Compact metric sets closed]\label{exr:ii06-expand-08}
Prove a compact subset of a metric space is closed.
\end{exercise}
\begin{hint}
For \(x\notin K\), use the positive minimum of \(d(x,\cdot)\) on \(K\) or a finite-ball argument.
\end{hint}
\begin{solution}
Distance from \(x\) to compact \(K\) attains a positive minimum; a smaller ball around \(x\) misses \(K\).
\end{solution}

\begin{exercise}[Heine--Cantor]\label{exr:ii06-expand-09}
Prove a continuous map on a compact metric domain is uniformly continuous.
\end{exercise}
\begin{hint}
Argue by sequences of increasingly close pairs and compact subsequences.
\end{hint}
\begin{solution}
A compact subsequence converges to one point on both sides; continuity contradicts any fixed output separation.
\end{solution}

\subsection*{Counterexamples and hypothesis testing}

\begin{exercise}[Closed bounded outside Euclidean finite dimension]\label{exr:ii06-expand-10}
Why can closed and bounded fail to imply compactness in infinite-dimensional normed spaces?
\end{exercise}
\begin{hint}
Look for an infinite uniformly separated sequence in the unit ball.
\end{hint}
\begin{solution}
Such a sequence has no Cauchy, hence no convergent, subsequence; the closed unit ball is then not compact.
\end{solution}

\begin{exercise}[Complete not compact]\label{exr:ii06-expand-11}
Give a complete metric space that is not compact.
\end{exercise}
\begin{hint}
Use \(\mathbb R\).
\end{hint}
\begin{solution}
\(\mathbb R\) is complete but unbounded, hence not compact.
\end{solution}

\begin{exercise}[Bounded not compact]\label{exr:ii06-expand-12}
Give a bounded noncompact subset of \(\mathbb R\).
\end{exercise}
\begin{hint}
Use an open interval.
\end{hint}
\begin{solution}
\((0,1)\) is bounded but not closed and not compact.
\end{solution}

\subsection*{Applications and investigations}

\begin{exercise}[Optimization guarantee]\label{exr:ii06-expand-13}
Why is compactness valuable before numerically maximizing a continuous objective?
\end{exercise}
\begin{hint}
Use the extreme-value theorem.
\end{hint}
\begin{solution}
It guarantees a true maximizer exists in the feasible set rather than only a supremum approached by escaping iterates.
\end{solution}

\begin{exercise}[Subsequence extraction]\label{exr:ii06-expand-14}
Why does compactness help prove existence from a sequence of approximations?
\end{exercise}
\begin{hint}
Extract a convergent subsequence and pass properties to the limit.
\end{hint}
\begin{solution}
It turns bounded/controlled approximations into a candidate limit inside the space.
\end{solution}

\subsection*{Challenge problems}

\begin{exercise}[Compact iff complete and totally bounded]\label{exr:ii06-expand-15}
Explain the difficult direction: complete and totally bounded implies compact in a metric space.
\end{exercise}
\begin{hint}
Use successively finer finite covers to construct a Cauchy subsequence from any sequence.
\end{hint}
\begin{solution}
Nested selections give a subsequence Cauchy at scales \(2^{-k}\); completeness gives a limit, hence sequential compactness and compactness.
\end{solution}

\begin{exercise}[Nested compact sets]\label{exr:ii06-expand-16}
Prove a decreasing sequence of nonempty compact sets \(K_n\) has nonempty intersection if all lie in one compact space.
\end{exercise}
\begin{hint}
Assume empty intersection and use complements as an open cover.
\end{hint}
\begin{solution}
The complements cover \(K_1\); a finite subcover implies some \(K_N\) is empty by nestedness, contradiction.
\end{solution}
% END VOL02-EXPANSION II06-exercises-01
\section*{Chapter summary}
The central definitions, estimates, and structural theorems of this chapter will be used repeatedly in the remaining analysis volume.
